% =====================================================================
% reports/preamble.tex - Preambolo condiviso dei report di milestone
% =====================================================================
% Contiene solo configurazione tipografica e pacchetti, nessun contenuto.
% Engine fissato a pdflatex in reports/.latexmkrc, per portabilità: i report
% devono compilare su una TinyTeX minimale senza font di sistema.
% Convenzioni ispirate a E:\harmony-book, adattate a un documento tecnico
% invece che editoriale: nessun font commerciale, nessuna classe memoir.
% ---------------------------------------------------------------------

\usepackage[T1]{fontenc}
\usepackage[utf8]{inputenc}
\usepackage{lmodern}
\usepackage[italian]{babel}
\usepackage[a4paper,margin=2.6cm]{geometry}
\usepackage{microtype}

% --- Tabelle e figure ---
\usepackage{booktabs}
\usepackage{tabularx}
\usepackage{longtable}
\usepackage{graphicx}
\usepackage[font=small,labelfont=bf]{caption}
\usepackage{xcolor}

% --- Matematica, per i conti su temporizzazione e capacità di canale ---
\usepackage{amsmath}

% --- Codice sorgente ---
\usepackage{listings}
\definecolor{grigiochiaro}{gray}{0.96}
\definecolor{commento}{rgb}{0.35,0.45,0.35}
\lstset{
  basicstyle=\ttfamily\footnotesize,
  backgroundcolor=\color{grigiochiaro},
  commentstyle=\color{commento}\itshape,
  keywordstyle=\bfseries,
  breaklines=true,
  showstringspaces=false,
  frame=single,
  framerule=0pt,
  xleftmargin=0.5em,
  captionpos=b,
}

% --- Riferimenti interni ---
\usepackage[hidelinks]{hyperref}

% --- Intestazioni ---
% Deliberatamente senza fancyhdr: non è presente nella TinyTeX di questa macchina e
% tlmgr rifiuta di installarlo senza aggiornare prima se stesso. Un report di milestone
% non ha bisogno di intestazioni personalizzate, e una dipendenza in meno vale più di
% una riga in testa alla pagina. Numerazione in basso, stile standard.
\pagestyle{plain}

% --- Unità di misura ---
% siunitx non è garantito presente su una TinyTeX minimale, e per un report tecnico
% servono solo grandezze semplici con lo spazio fine prima dell'unità. Si definisce
% quindi la macro solo se il pacchetto non l'ha già fornita, così che il corpo del
% documento resti valido anche dove siunitx c'è.
\providecommand{\SI}[2]{#1\,#2}

% --- Macro del progetto ---
% Un byte in esadecimale, scritto sempre nello stesso modo.
\newcommand{\hx}[1]{\texttt{0x#1}}
% Una fonte di livello n della gerarchia dichiarata in SOURCES.md.
\newcommand{\liv}[1]{\textsuperscript{L#1}}
